In the construction every one-loop component of a local operator along a generator $x$ of $\mathfrak{psl}(4|4)$ is read from the defect flow of Lemma \ref{lem:c8-transfer}: it is the sum over the pairs of letters of the coefficients of $\varepsilon^0$ of the defect cycles, cycles through the two letters with one special edge and one vertex $S^{(1)}(x)$, minus the bracket of the operator with the arity-one finite counterterm $\hbar F_1^{(1)}(x)$, which solves the arity-one equation of the action with all its free directions. This master formula (Theorem \ref{thm:c11-classes}) needs an argument: the graphs with the special edge on a bridge have finite limits, which cancel exactly against the special-edge parts of their counterterm partners. The chapter first treats the composite term of Chapter 10, the second evaluation of $D_1$ and $P_1$ through the scaling of the regulator, in which $S^{(1)}(\Delta)$ or $S^{(1)}(P_b)$ sits on the cycle as a line from a ghost leg to a gauge leg (Theorem \ref{thm:c10-D1}, Proposition \ref{prop:c10-P1}); on the letters $A|_V$, $\lambda$, $c$ the renormalised composite two-point functional contains $\log|p|^2$ of the momentum (Remark \ref{rem:c9-Pi}), and the locality of the composite term is proved: $Q$ maps the composite graphs of a cocycle to cycles with one special edge, and the composite term is a sum over the pairs of letters of the composite pair kernels. These kernels are simplex integrals, evaluated on twist two, the $\mathfrak{su}(1|1)$ sector, the fermion and gluon $\mathfrak{sl}(2)$ sectors and two mixed families; with the sector coefficients $-\tfrac13$ units of $\hbar F_1^{(1)}(\Delta)$ computed from the action-level arity-two equations (Theorem \ref{thm:c9-theta2}), $D_1$ is four times the map of Theorem \ref{thm:c10-brst}, with spectrum twice the one-loop anomalous dimensions, and $P_1$ vanishes on cohomology in the translation-covariant scheme; the operator-level arity-two equation reproduces the value $-\tfrac13$ family by family. The last section derives the master formula along every generator, computes the action-level obstructions and solutions along the odd generators and along $K_b$ with their free directions and the arity-two equations that relate them, assembles $Q_1$, $S_1$, $D_1$, $P_1$ and $K_1$ on the scalar bilinears as representatives together with their scheme-independent content, and compares with the one-loop differential of the holomorphic twist at the level of classes. Chapter 12 states the construction for every generator and every arity together (Theorems \ref{thm:c12-master}, \ref{thm:c12-ar2} and \ref{thm:c12-F0}).
